% Nouveau EGA V (cote 25): archive cote 25, site page 41 (PDF page 42).
% typescript. Draft compared: none.
% AI-transcribed by Claude Opus 5.5 (claude-opus-5-5, adaptive thinking, effort high), Message Batches API, batch ega25b (msgbatch_01A7n5YuXubUPNK1zsepi8RG), 2026-11-13; not yet checked by eye.
\scanpage{25/41}
% NOTE: the first typed line at the head of the page ("morphisme de préschémas lisses sur S de dim relative r-1 et r:") is enclosed by hand in a loop with an insertion hook; the same hook stands at the right margin after "du" in b). It is transcribed there, as \inserted[typed]{...}.
\mspage{\unsure{37}}

passab\unsure{le}ment de sueur déjà), on aurait pu remarquer que, \var{\deleted{Comme}}{}{deletion}{ega25_25_041-v0001} de
toutes façons du point de vue ensembliste $\underline{Y}^{sing}$ \var{\inserted[typed]{est de}}{est de}{insertion}{ega25_25_041-v0002} \var{\xout{\unsure{des fibres de}}}{}{typewriter strike-out}{ega25_25_041-v0003}dimen%
sion r-1 \var{\xout{\unsure{sur S, alors que si}}}{}{typewriter strike-out}{ega25_25_041-v0004} si S=Spec(k), alors que $\overset{\hand{\vee}}{P}$ est de dimen-%
sion r, \var{\inserted{donc}}{donc}{insertion}{ega25_25_041-v0005} l'image de $\underline{Y}^{sing}$ dans $\overset{\hand{\vee}}{P}$ est de codimension $\unsure{\geq}$1, ce qui redonne
2.12. (l'argument bien entendu n'est pas essentiellement distinct de
celui utilisé dans 2.12.). On notera que le plus souvent, cet ensemble
est effectivement de codimension 1 (cf plus bas). Par suite, on ne
pourra trouver en général des ``pinceaux linéaires'' de sections hyperpla\var{\xout{n}}{}{typewriter strike-out}{ega25_25_041-v0006}%
nes toutes \var{\inserted[typed]{\var{\xout{\unsure{géom}}}{}{typewriter strike-out}{ega25_25_041-v0008}}}{\var{\xout{\unsure{géom}}}{}{typewriter strike-out}{ega25_25_041-v0008}}{insertion}{ega25_25_041-v0007} \var{\xout{\unsure{non singulières}}}{}{typewriter strike-out}{ega25_25_041-v0009} lisses. Cependant, nous allons voir qu'on
pourra s'arranger souvent pour \var{\xout{que}}{}{typewriter strike-out}{ega25_25_041-v0010} trouver des pinceaux formés de sec-%
tions hyperplanes n'ayant aucun point supersingulier, grâce au fait que
dans les cas les plus fréquents, l'image de $\underline{Y}^{sup\ sing}$ dans $\overset{\hand{\vee}}{P}$ sera de
codimension 2. Nous allons d'abords résumer les points essentiels
de nature différentielle de la situation étudiée ici:

\uline{Théorème} 8.\var{\over{\illegible{1}}{7}}{7}{writing over}{ega25_25_041-v0011}. a) Le sous-préschéma $\underline{Y}^{sing}$ (défini au par.16) dans
la situation actuelle n'est autre que $P(\underline{\hand{\unsure{\nu}}}_{X/P})$ , considéré comme
sous-préschéma de $\underline{Y}$ comme expliqué plus \var{\xout{h}}{}{typewriter strike-out}{ega25_25_041-v0012}haut.

b) L'ensemble sous-jacent au préschéma $\underline{Y}^{sup\ sing}$ (cf par. 16), \var{\xout{\illegible{1}}}{}{typewriter strike-out}{ega25_25_041-v0013}
\var{\deleted{\unsure{fermé de}}}{}{deletion}{ega25_25_041-v0014} n'est autre que l'ensemble des points de ramification du \var{\inserted[typed]{morphisme de préschémas lisses sur S de dim relative r-1 et r:}}{morphisme de préschémas lisses sur S de dim relative r-1 et r:}{insertion}{ega25_25_041-v0015}
$\underline{Y}^{sing} = P(\underline{\hand{\unsure{\nu}}}_{X/P}) \longrightarrow \check{P}$ , i.e. pour que ce dernier morphisme soit
non ramfié en le point \var{\over{\illegible{1}}{$y$}}{$y$}{writing over}{ega25_25_041-v0016} (réf à définition) il f. et s. que \var{\over{\illegible{1}}{$y$}}{$y$}{writing over}{ega25_25_041-v0017}
soit géométriquement un point singulier ordinaire pour $\phi_{\hand{\xi}}$ ($\hand{\xi}$ étant le
point de $\overset{\hand{\vee}}{P}$ image de \var{\over{\illegible{1}}{$y$}}{$y$}{writing over}{ega25_25_041-v0018} ).

c) Supposons S=Spec(k), et que \var{\over{\illegible{1}}{$y$}}{$y$}{writing over}{ega25_25_041-v0019} \var{\inserted{$\in$}}{$\in$}{insertion}{ega25_25_041-v0020}\var{\inserted[typed]{$\underline{Y}^{sing} = P(\underline{\hand{\unsure{\nu}}})$}}{$\underline{Y}^{sing} = P(\underline{\hand{\unsure{\nu}}})$}{insertion}{ega25_25_041-v0021} soit rationnel sur k; soient x
et $\hand{\xi}$ ses projections dans X(k) re\unsure{sp|pp} $\overset{\hand{\vee}}{P}$(k) , et considérons la \var{\xout{\unsure{sous-espace}}}{}{typewriter strike-out}{ega25_25_041-v0022}
sous-variété linéaire \var{\inserted{$H'$}}{$H'$}{insertion}{ega25_25_041-v0023} de $\overset{\hand{\vee}}{P}$ ``image'' de l'application tangente de

\var{\margin{41}}{\placedmargin{41}}{marginal note placed}{ega25_25_041-v0024}
